\documentclass[twocolumn,prc,amsmath,amssymb,superscriptaddress]{revtex4-2}

\usepackage{graphicx}
\usepackage{hyperref}
\usepackage{booktabs}
\usepackage{xcolor}
\usepackage{listings}

\lstset{
  language=Python,
  basicstyle=\ttfamily\footnotesize,
  keywordstyle=\color{blue},
  stringstyle=\color{red!60!black},
  commentstyle=\color{green!50!black},
  breaklines=true,
  frame=single,
  xleftmargin=1em,
  xrightmargin=0.5em,
}

\begin{document}

\title{Nuclear Shell Model Hamiltonians on the Hyperspherical Lattice}

\author{J.~Loutey}
\affiliation{Independent Researcher, Kent, Washington}

\date{April 9, 2026}

\begin{abstract}
We present structurally sparse qubit Hamiltonians for the nuclear shell
model, constructed by applying the GeoVac angular-sparsity framework to
the harmonic oscillator (HO) single-particle basis with a Mayer--Jensen
spin--orbit term. The HO shell closures $2, 8, 20, 40, 70, 112$ emerge
directly from graph state counting, and the physical magic numbers
$2, 8, 20, 28, 50, 82, 126$ are recovered with a spin--orbit strength
$v_{ls}/\hbar\omega \approx 0.17$ and a Nilsson $\ell(\ell+1)$ correction
$d_{\ell\ell}/\hbar\omega \approx 0.02$. We build the first two-species
qubit encoding for the deuteron ($1p+1n$, $16$ qubits, $592$ non-identity
Pauli terms, $1$-norm $\approx 227$~MeV) using the Minnesota
nucleon--nucleon potential and Moshinsky--Talmi brackets, and extend the
architecture to helium-4 ($2p+2n$, $16$ qubits, $712$ Pauli terms) with
intra-species antisymmetrization and proton--proton Coulomb repulsion.
Despite a Hilbert-space dimension $12.25\times$ larger than the deuteron,
He-4 requires only $1.20\times$ more Pauli terms --- a direct consequence
of the angular selection rules that carry through to multi-particle
systems. A composed nuclear--electronic deuterium encoding ($26$ qubits)
validates the multi-block architecture and quantifies the
$\sim 10^{13}$ coefficient ratio between nuclear, electronic, and
hyperfine scales. We state and apply the Fock projection rigidity
theorem: the $S^3$ conformal projection that makes electronic Hamiltonians
equivalent to a graph Laplacian is unique to the Coulomb potential and
does \emph{not} extend to HO or Woods--Saxon systems. For nuclear
Hamiltonians, the GeoVac contribution is therefore the angular machinery
and qubit encoding, not the conformal projection. All nuclear physics
inputs (Minnesota parameters, spin--orbit strength, $\hbar\omega$)
follow the standard shell-model literature; the quantitative binding
energies at $N_{\mathrm{shells}}=2$ are far from experiment and should
be read as encoding-validation benchmarks, not as a nuclear-structure
calculation.
\end{abstract}

\maketitle

%% =====================================================================
\section{Introduction}
%% =====================================================================

Nuclear quantum simulation is an active frontier in quantum
computing~\cite{Roggero2020,Hagen2014,Romero2022}. Exact solvers
(no-core shell model, NCSM~\cite{Navratil2000,BarrettNavratilVary2013})
remain the gold standard for light nuclei $A \lesssim 16$, but scale
factorially with the number of active nucleons, and their extension to
medium-mass nuclei requires importance truncation or symmetry-adapted
bases. On current and near-term quantum hardware, the practical
bottleneck is the resource cost of the qubit Hamiltonian: the nuclear
interaction is short-ranged and contains tensor and three-body
components that produce dense Pauli representations in conventional
encodings. Resource estimates for ab initio nuclear simulation using
second-quantized Hamiltonians therefore sit at the edge of current
hardware feasibility.

The GeoVac framework~\cite{Paper14,Paper22} exploits angular-momentum
selection rules to produce structurally sparse qubit Hamiltonians for
molecular electronic structure, with Pauli counts that are orders of
magnitude below Gaussian-basis baselines for the same system size. The
angular selection rules, which are derived from Wigner $3j$ symbols and
the Gaunt integral, are \emph{independent of the radial potential}:
they depend only on the single-particle angular basis, not on whether
the interaction is Coulombic, Yukawa, or nucleon--nucleon. This paper
investigates whether that sparsity carries over to nuclear shell model
Hamiltonians.

We report three results. First, the harmonic oscillator shell closures
and the Mayer--Jensen magic numbers are reproduced on the GeoVac
$(n_r, \ell, m_\ell, \sigma)$ graph via an explicit spin--orbit term.
Second, the deuteron and He-4 qubit Hamiltonians, built from the
Minnesota NN potential using Moshinsky--Talmi brackets, exhibit Pauli
counts that grow sub-linearly in the Hilbert-space dimension --- the
deuteron has $16$ qubits and $592$ Pauli terms, while He-4 has the
same $16$ qubits but only $712$ Pauli terms despite a $12.25\times$
larger Hilbert space. Third, a composed nuclear--electronic encoding
for deuterium demonstrates that the block architecture used for
core--valence molecular systems extends to nucleon--electron systems,
though the $10^{13}$ coefficient ratio between nuclear-scale and
hyperfine terms makes a single-pass quantum simulation impractical
without block partitioning.

We also state the Fock projection rigidity theorem, which identifies
the parts of the GeoVac framework that transfer to non-Coulomb systems
(the angular machinery) and the parts that do not (the $S^3$ conformal
projection). This delineation is the main conceptual contribution of
the paper.

The paper is organized as follows. Section~\ref{sec:shell} constructs
the nuclear shell model on the GeoVac lattice. Section~\ref{sec:rigidity}
states the Fock projection rigidity theorem. Sections~\ref{sec:deuteron}
and~\ref{sec:he4} present the deuteron and He-4 qubit Hamiltonians.
Section~\ref{sec:ne} presents the composed nuclear--electronic
encoding. Sections~\ref{sec:discussion} and~\ref{sec:conclusion}
conclude.

%% =====================================================================
\section{Nuclear Shell Model on the Hyperspherical Lattice}
\label{sec:shell}
%% =====================================================================

\subsection{Harmonic oscillator basis and shell closures}

The standard starting point for the nuclear shell model is the
three-dimensional isotropic harmonic oscillator with single-particle
energies
\begin{equation}
  E_N^{\mathrm{HO}} = \hbar\omega \left(N + \tfrac{3}{2}\right),
  \qquad
  N = 2 n_r + \ell,
\label{eq:ho-energy}
\end{equation}
and degeneracy $(N+1)(N+2)$ per major shell including spin. Summing
these degeneracies yields the HO shell closures
\[
  2, \; 8, \; 20, \; 40, \; 70, \; 112, \; \ldots
\]
These are the correct magic numbers of a pure 3D isotropic oscillator
but disagree with the experimentally observed nuclear magic numbers
above $N=20$.

On the GeoVac lattice, each state $(n_r, \ell, m_\ell, \sigma)$ is a
graph node with diagonal energy given by
Eq.~\eqref{eq:ho-energy}. Because the HO shell spacing is uniform in
$N$, the HO magic numbers are recovered by simple state counting at any
choice of kinetic coupling; they do not depend on the graph off-diagonal
structure. This is a numerical confirmation of the packing counting
rule from Paper~0~\cite{Paper0} as applied to the HO shell degeneracies.

\subsection{Spin--orbit term and the real magic numbers}

Mayer and Jensen~\cite{Mayer1949,Jensen1949} recognized that the
observed nuclear magic numbers $28, 50, 82, 126$ require an additional
spin--orbit interaction whose sign is \emph{opposite} to the atomic
case: the nuclear spin--orbit coupling lowers the $j = \ell + 1/2$
sub-shell below $j = \ell - 1/2$. Adding the Nilsson $\ell(\ell+1)$
term, which mimics the flattening of the Woods--Saxon mean field
relative to the pure HO, one writes
\begin{equation}
  E(N,\ell,j) = \hbar\omega\left(N + \tfrac{3}{2}\right)
  - v_{ls}\,\langle\boldsymbol{\ell}\cdot\mathbf{s}\rangle
  - d_{\ell\ell}\,\ell(\ell+1),
\label{eq:mayer-jensen}
\end{equation}
with $\langle \boldsymbol{\ell}\cdot\mathbf{s}\rangle = \ell/2$ for
$j = \ell + 1/2$ and $-(\ell+1)/2$ for $j = \ell - 1/2$.

A two-parameter scan over $(v_{ls}, d_{\ell\ell})$ identifies the range
of parameters for which all seven physical magic numbers $2, 8, 20, 28,
50, 82, 126$ are recovered simultaneously as gaps in the single-particle
spectrum. With seven HO shells and the midpoint of the valid region we
obtain
\begin{equation}
  \frac{v_{ls}}{\hbar\omega} \approx 0.171, \qquad
  \frac{d_{\ell\ell}}{\hbar\omega} \approx 0.021 ,
\label{eq:optimal}
\end{equation}
which is consistent with the Mayer--Jensen parameters tabulated in
standard shell-model references~\cite{Suhonen2007,Talmi1993}. These
are nuclear-physics inputs; the GeoVac framework contributes nothing
to their determination. Table~\ref{tab:levels} gives the level
ordering, with shell closures marked.

\begin{table}[t]
  \centering
  \caption{Level ordering from Eq.~\eqref{eq:mayer-jensen} at
    $v_{ls}/\hbar\omega = 0.171$, $d_{\ell\ell}/\hbar\omega = 0.021$.
    Shell closures at the experimentally observed magic numbers are
    marked with $\star$.}
  \label{tab:levels}
  \begin{ruledtabular}
  \begin{tabular}{llr}
    Shell & Levels (in order) & Cumul. \\
    \hline
    1 & $0s_{1/2}$ & $2^\star$ \\
    2 & $0p_{3/2},\;0p_{1/2}$ & $8^\star$ \\
    3 & $0d_{5/2},\;1s_{1/2},\;0d_{3/2}$ & $20^\star$ \\
    4 & $0f_{7/2}$ & $28^\star$ \\
    5 & $1p_{3/2},\;0f_{5/2},\;1p_{1/2},\;0g_{9/2}$ & $50^\star$ \\
    6 & $1d_{5/2},\;0h_{11/2},\;2s_{1/2},\;0g_{7/2},\;1d_{3/2}$ & $82^\star$ \\
    7 & $1f_{7/2},\;0i_{13/2},\;2p_{3/2},\;0h_{9/2},\;1f_{5/2},\;2p_{1/2}$
      & $126^\star$ \\
  \end{tabular}
  \end{ruledtabular}
\end{table}

The spin--orbit operator is implemented directly in the
$(m_\ell, m_s)$ basis via
\[
  \boldsymbol{\ell}\cdot\mathbf{s} = \ell_z s_z
  + \tfrac{1}{2}(\ell_+ s_- + \ell_- s_+),
\]
producing an off-diagonal coupling between $(m_\ell, {\uparrow})$ and
$(m_\ell + 1, {\downarrow})$ within each $(n_r, \ell)$ subspace. The
two-body electron-repulsion-like integrals are unchanged because
spin--orbit is a one-body operator. The net effect on single-particle
Hamiltonian sparsity is a modest increase from pure diagonal to a
block-tridiagonal structure within each $(n_r, \ell)$ submatrix; the
overall qubit encoding cost is dominated by the two-body terms in all
cases we consider.

%% =====================================================================
\section{The Fock Projection Rigidity Theorem}
\label{sec:rigidity}
%% =====================================================================

The GeoVac framework for electronic structure uses a conformal
projection, originally due to Fock~\cite{Fock1935}, that maps the
bound-state spectrum of the hydrogen atom onto the free Laplacian on
the three-sphere $S^3$. The projection is
\[
  p_0 = \sqrt{-2 E_n} = Z/n ,
\]
and it works because the hydrogen spectrum is $\ell$-degenerate within
each principal quantum number $n$: all orbitals $(n,\ell,m)$ with the
same $n$ have the same energy. The conformal projection then identifies
the radial Schr\"odinger equation on $\mathbb{R}^3$ with the Laplace
equation on $S^3$, and the graph Laplacian on the $(n,\ell,m)$ lattice
converges to the continuous $S^3$ Laplacian in the asymptotic limit
(Paper~7~\cite{Paper7}).

For nuclear shell model Hamiltonians the projection fails, for reasons
that are elementary but worth stating precisely. We summarize the
result as

\textbf{Theorem (Fock projection rigidity).}
\emph{The $S^3$ Fock projection $p_0 = \sqrt{-2 E_n}$ maps a one-particle
central-field Hamiltonian onto the free Laplacian on $S^3$ if and only
if the spectrum $E_{n\ell}$ is independent of $\ell$ within each $n$.
The Coulomb potential $-Z/r$ is the unique central potential exhibiting
this accidental degeneracy, as a consequence of its hidden SO(4)
symmetry (the Runge--Lenz vector). Any deformation
$V(r) = -Z/r + \delta V(r)$ with $\delta V \not\equiv 0$ lifts the
$\ell$-degeneracy, breaking the conformal equivalence between the
radial Schr\"odinger problem and the $S^3$ graph Laplacian.}

The proof is standard and amounts to observing that Fock's
stereographic construction in momentum space requires the hydrogenic
relation $E_n \propto -1/n^2$ with no $\ell$-dependence; this fails for
any potential that does not have the SO(4) symmetry of the pure Coulomb
problem~\cite{Bander1966}. The consequence for the GeoVac lattice is
that the graph node weights, which depend only on $n$ in the Coulomb
case, cannot absorb an $\ell$-dependent correction without reducing the
graph Laplacian to the trivial diagonal of the corrected spectrum. In
that regime the off-diagonal adjacency structure contributes nothing.

The relevant central potentials for nuclear physics --- the 3D
harmonic oscillator, the Woods--Saxon well, the square well, and the
Yukawa potential --- all have $\ell$-dependent spectra. The 3D harmonic
oscillator has its own accidental degeneracy (the SU(3) group of the
isotropic oscillator groups states with common $N = 2n_r + \ell$),
but the degeneracy pattern is different from Coulomb's SO(4), and no
stereographic projection onto a compact manifold produces the graph
Laplacian form that underlies the GeoVac electronic framework. The
Bargmann--Segal transform~\cite{Bargmann1961} is the natural analogue
for the HO, mapping the $L^2$ Hilbert space to a holomorphic
reproducing-kernel space on $\mathbb{C}^3$; whether this can be
discretized into a graph whose Laplacian converges to the HO
Hamiltonian is an open question.

\emph{Negative result, positive implication.} The rigidity theorem
tells us what \emph{cannot} be carried over from Papers~7 and 11--15
to the nuclear setting: the conformal projection, the claim that
the $S^3$ graph Laplacian \emph{is} the quantum kinetic operator in
some dual description, and the $-1/16$ kinetic constant. What
\emph{can} be carried over is the angular machinery: the Wigner $3j$
coupling coefficients, the Gaunt integral selection rules, the block
structure of the two-body interaction tensor, and the O($Q^{2.5}$)
Pauli scaling that follows~\cite{Paper22}. For nuclear shell models,
we use the angular machinery to assemble the qubit Hamiltonian and use
explicit radial matrix elements (Slater-like integrals computed from
HO wavefunctions) for the radial part of the two-body interaction. The
radial part is no longer related to the graph; it is a precomputed
numerical table.

%% =====================================================================
\section{The Deuteron: First Nuclear Qubit Hamiltonian}
\label{sec:deuteron}
%% =====================================================================

\subsection{System and interaction}

The deuteron is a bound state of one proton and one neutron in the
$^3S_1$ channel, with experimental binding energy $2.2246$ MeV and an
rms charge radius of $\sim 2.1$ fm. The large radius is a consequence
of the small binding; the deuteron wavefunction extends well beyond
the range of the nuclear force, and accurate binding-energy
calculations require oscillator bases with many major shells. We use
the Minnesota potential~\cite{Thompson1977}, a phenomenological
soft-core central interaction composed of three Gaussian terms:
\begin{equation}
  V_{\mathrm{NN}}(r) = \sum_{i=1}^{3} V_i \, e^{-\kappa_i r^2}
  \left[\tfrac{1}{2}(u + (2 - u) P^\sigma)\right],
\label{eq:minnesota}
\end{equation}
with $(V_R, V_S, V_T) = (200, -178, -91.4)$ MeV and the standard
$u$ and range parameters of Ref.~\cite{Thompson1977}. The operator
$P^\sigma$ is the spin exchange. The two-body matrix elements in the
HO basis are computed by transforming the antisymmetrized
two-nucleon wavefunction to relative and center-of-mass coordinates
using Moshinsky--Talmi brackets~\cite{Moshinsky1959,Talmi1952}.

\subsection{Two-species qubit encoding}

Protons and neutrons are distinct fermion species for the purposes of
antisymmetrization: Pauli exclusion applies within each species but
not between them. We therefore allocate two separate Jordan--Wigner
registers, one of size $Q_p$ for the proton orbitals and one of size
$Q_n$ for the neutron orbitals, and combine them by tensor product.
At $N_{\mathrm{shells}} = 2$ (major shells $N=0$ and $N=1$) each
species has
\[
  Q_p = Q_n =
  \underbrace{2}_{0s_{1/2}} + \underbrace{6}_{0p} = 8,
\]
giving a total $Q = 16$ qubits. The single-particle states are labeled
in the uncoupled $(n_r, \ell, m_\ell, m_s)$ basis; the coupling to
$j$-scheme is performed implicitly via the Minnesota matrix elements.

The one-body Hamiltonian is
\[
  h_1 = \hbar\omega (2 n_r + \ell + \tfrac{3}{2}) \, \delta_{ij},
\]
diagonal in the HO basis. The two-body Hamiltonian is
\[
  \hat V_{\mathrm{pn}} = \sum_{ijkl} V_{ijkl}
  \, a^\dagger_{p,i} a^\dagger_{n,j} a_{n,l} a_{p,k},
\]
where $V_{ijkl}$ is computed from Eq.~\eqref{eq:minnesota} via
Moshinsky--Talmi brackets. Because protons and neutrons are on
different registers, each term factorizes as
\[
  a^\dagger_{p,i} a_{p,k} \otimes a^\dagger_{n,j} a_{n,l},
\]
yielding a tensor-product Pauli string that has support on both
registers.

\subsection{Resource counts}

Table~\ref{tab:deuteron} summarizes the qubit Hamiltonian at
$N_{\mathrm{shells}}=2$, $\hbar\omega = 10$ MeV. The non-identity Pauli
term count is $592$, with a $1$-norm of $\approx 227$ MeV dominated by
the HO diagonal and the Minnesota spin-triplet contribution. Of the
$592$ terms, $80$ are $Z$-only (number-operator products) and the
remaining $512$ carry at least one $X$ or $Y$, reflecting the
off-diagonal angular couplings generated by the Minnesota operator.

\begin{table}[t]
  \centering
  \caption{Deuteron qubit Hamiltonian at $N_{\mathrm{shells}}=2$,
    $\hbar\omega = 10$ MeV. ``Non-$I$'' denotes the non-identity Pauli
    terms. Hilbert-space dimension is
    $\binom{Q_p}{1}\binom{Q_n}{1} = 64$.}
  \label{tab:deuteron}
  \begin{ruledtabular}
  \begin{tabular}{lr}
    Quantity & Value \\
    \hline
    $N_{\mathrm{shells}}$ & 2 \\
    Qubits $Q = Q_p + Q_n$ & $8 + 8 = 16$ \\
    Non-$I$ Pauli terms & $592$ \\
    $Z$-only terms & $80$ \\
    $XY$-containing terms & $512$ \\
    $1$-norm (non-$I$) & $\approx 227.3$ MeV \\
    Hilbert dim. ($1p+1n$) & $64$ \\
  \end{tabular}
  \end{ruledtabular}
\end{table}

\subsection{Binding energy at $N_{\mathrm{shells}}=2$}

Diagonalizing the $(Q_p = 8)\times(Q_n = 8) = 64$-dimensional FCI
matrix at $\hbar\omega = 10$ MeV and subtracting the HO zero-point
energy gives a bound-state energy that is far from the experimental
$-2.2246$ MeV. This is expected: the deuteron's rms radius of
$\sim 2.1$ fm is comparable to the HO length at $\hbar\omega = 10$
MeV, so $N_{\mathrm{shells}}=2$ has insufficient radial flexibility to
describe the long tail of the wavefunction. Converged Minnesota
calculations require $N_{\mathrm{shells}} \gtrsim 6$--$8$. We therefore
emphasize that the results in Table~\ref{tab:deuteron} are
\emph{encoding-validation} benchmarks, not nuclear-structure
predictions. The Pauli count and $1$-norm characterize the cost of
the qubit Hamiltonian, and the FCI diagonalization verifies that the
Pauli encoding reproduces the matrix representation.

%% =====================================================================
\section{Helium-4: Antisymmetrization Within Species}
\label{sec:he4}
%% =====================================================================

We extend the deuteron encoding to He-4, a doubly magic nucleus with
two protons and two neutrons, experimental binding energy $28.296$
MeV. The architecture is the same as for the deuteron, with two
changes:

\begin{enumerate}
  \item Intra-species antisymmetrization: protons fill a $2$-particle
    Slater determinant on the proton register, neutrons fill a
    $2$-particle Slater determinant on the neutron register. The
    resulting two-body terms include $pp$ and $nn$ antisymmetrized
    matrix elements in addition to the $pn$ cross-species terms.
  \item Coulomb repulsion between the two protons, added as a
    conventional two-body Coulomb matrix element in the HO basis
    (approximately $0.7$ MeV for the ground state).
\end{enumerate}

Both changes leave the number of qubits unchanged ($Q = 16$) because
adding particles does not add spin orbitals; they only increase the
Hilbert-space dimension. At $N_{\mathrm{shells}}=2$ the $pp$ pair
occupies a $2$-particle subspace of dimension $\binom{8}{2} = 28$, and
similarly for $nn$, so the full He-4 Hilbert space has dimension
$28 \times 28 = 784$.

Table~\ref{tab:he4} compares He-4 to the deuteron. The key observation
is that the Pauli count grows by only $20\%$ ($592 \rightarrow 712$)
even though the Hilbert-space dimension grows by a factor of
$784/64 = 12.25$. The angular selection rules responsible for the
deuteron sparsity carry over to the antisymmetrized pairs without
introducing new classes of terms; only the $pp$ and $nn$ blocks are
added, and those blocks share the same angular structure as the $pn$
block. The Coulomb correction is a small additional contribution that
does not add any new Pauli terms (it modifies only the existing $pp$
coefficients).

\begin{table}[t]
  \centering
  \caption{Comparison of deuteron and He-4 qubit Hamiltonians at
    $N_{\mathrm{shells}}=2$. Both systems use $16$ qubits; the
    difference is the Hilbert-space dimension and the number of
    particles. Pauli counts refer to non-identity terms.}
  \label{tab:he4}
  \begin{ruledtabular}
  \begin{tabular}{lrrr}
    System & $Q$ & Pauli & Hilbert dim. \\
    \hline
    Deuteron ($1p+1n$)       & 16 & 592 & 64 \\
    He-4 ($2p+2n$, no Coulomb) & 16 & 712 & 784 \\
    He-4 ($2p+2n$, + Coulomb)  & 16 & 712 & 784 \\
  \end{tabular}
  \end{ruledtabular}
\end{table}

The $1$-norm of the He-4 Hamiltonian at $\hbar\omega = 10$ MeV is
$\approx 557$ MeV (without Coulomb) and $\approx 552$ MeV (with
Coulomb), reflecting the dominance of the HO diagonal piece. The
ground-state energy from FCI diagonalization is far from the
experimental $-28.296$ MeV at $N_{\mathrm{shells}}=2$, for the same
reason as in the deuteron: the HO basis is too small to describe the
$\alpha$-particle charge distribution. Converged Minnesota
calculations of He-4 require $N_{\mathrm{shells}} \gtrsim 6$. As with
the deuteron, we emphasize that Table~\ref{tab:he4} is an
encoding-validation benchmark. The physical result of interest is the
sub-linear Pauli scaling, which is a direct consequence of the
angular-sparsity theorem of Paper~22~\cite{Paper22} and confirms that
the theorem applies to multi-particle systems beyond the single-nucleon
case.

%% =====================================================================
\section{Composed Nuclear--Electronic Deuterium}
\label{sec:ne}
%% =====================================================================

The composed architecture used for core--valence molecular electronic
structure~\cite{Paper14,Paper17} extends naturally to the
nuclear--electronic case: the nuclear block is the deuteron
Hamiltonian from Sec.~\ref{sec:deuteron}, the electronic block is a
single-electron hydrogen Hamiltonian on the $(n, \ell, m)$ lattice at
$n_{\max}=2$, and the coupling between them is given by two
one-body-like operators.

\subsection{Blocks and couplings}

\textbf{Nuclear block (Block A).} One proton and one neutron in the
HO basis at $N_{\mathrm{shells}}=2$, $16$ qubits, as described in
Sec.~\ref{sec:deuteron}. Energy scale $\sim 10^1$--$10^2$ MeV after
conversion to atomic units.

\textbf{Electronic block (Block B).} A single electron in the
hydrogenic $1s$--$2p$ manifold (Coulomb potential $-1/r$), $Q_e = 10$
qubits (5 spatial orbitals $\times$ 2 spins: $1s$, $2s$, $2p_{-1}$,
$2p_0$, $2p_{+1}$). The ground state is the exact hydrogenic $1s$ at
$-0.5$ Ha.

\textbf{Finite-size coupling.} A one-body operator on the electronic
register shifts the $1s$ energy by the finite nuclear-size correction
$\Delta E_{1s} = (2/5) Z^4 R_{\mathrm{nuc}}^2 / n^3$
\cite{Foldy1958,Friar1979}. At the proton charge radius
$R_{\mathrm{nuc}} = 0.8414$ fm, the shift is
$\Delta E \approx 1.01 \times 10^{-10}$ Ha, far below the subtraction
precision of float$64$ arithmetic when combined with the nuclear-scale
terms. The shift is reported analytically rather than by numerical
subtraction.

\textbf{Hyperfine coupling.} The proton spin $\mathbf{I}$ couples to
the electron spin $\mathbf{S}$ via
$H_{\mathrm{hf}} = A_{\mathrm{hf}} \, \mathbf{I}\cdot\mathbf{S}$, where
the hyperfine constant for the hydrogen $21$-cm transition is
$A_{\mathrm{hf}} = 1420.4$ MHz $\approx 2.16 \times 10^{-7}$ Ha. The
singlet--triplet gap is $A_{\mathrm{hf}}$, and the maximum shift of
any single ground-state manifold member is $3 A_{\mathrm{hf}}/4
\approx 1.62 \times 10^{-7}$ Ha.

The composed Hamiltonian is
\begin{equation}
  H = H_{\mathrm{nuc}} \otimes I_{e}
    + I_{\mathrm{nuc}} \otimes H_{e}
    + V_{\mathrm{fs}} + H_{\mathrm{hf}},
\label{eq:composed-ne}
\end{equation}
with $V_{\mathrm{fs}}$ the finite-size shift on the electronic register
and $H_{\mathrm{hf}}$ the cross-register hyperfine operator.

\subsection{Resource counts and coefficient hierarchy}

The composed Hamiltonian uses $Q = 16 + 10 = 26$ qubits. The Pauli
count on the full register is $614$ non-identity terms, of which
$592$ come from the nuclear block, $10$ from the electronic block
diagonal (one $Z$ per spin-orbital plus the absorbed finite-size
shift), and $12$ from the cross-register hyperfine operator. The
finite-size shift is absorbed into the diagonal electronic terms and
does not add new Pauli strings.

The coefficient hierarchy spans $\sim 13$ orders of magnitude
(Table~\ref{tab:ne-scales}): the nuclear diagonal dominates at
$\sim 3 \times 10^5$ Ha (after conversion from MeV to
Ha via $1$ Ha $\approx 27.2$ eV), the electronic binding is $\sim 0.5$ Ha,
the hyperfine coefficient is $\sim 10^{-8}$ Ha per Pauli string (the
$A_{\mathrm{hf}}/16$ from the $\mathbf{I}\cdot\mathbf{S}$ decomposition),
and the finite-size shift is $\sim 10^{-10}$ Ha but does not appear in
the Pauli list because it is below the float$64$ subtraction floor
relative to the nuclear-scale diagonal.

\begin{table}[t]
  \centering
  \caption{Coefficient hierarchy in the composed nuclear--electronic
    deuterium Hamiltonian (atomic units).}
  \label{tab:ne-scales}
  \begin{ruledtabular}
  \begin{tabular}{lc}
    Contribution & Scale (Ha) \\
    \hline
    Max Pauli coeff (nuclear)    & $\sim 3 \times 10^{5}$ \\
    Electronic binding scale     & $\sim 10^{-1}$ \\
    Hyperfine coeff (per Pauli)  & $\sim 10^{-8}$ \\
    Finite-size shift (analytic) & $\sim 10^{-10}$ \\
    \hline
    Pauli coeff ratio max/min    & $\sim 2 \times 10^{13}$ \\
  \end{tabular}
  \end{ruledtabular}
\end{table}

\subsection{Multi-scale precision problem}

The $\sim 10^{13}$ coefficient ratio is the central obstacle to
running Eq.~\eqref{eq:composed-ne} as a single-pass quantum simulation
on near-term hardware. To resolve a hyperfine shift of $10^{-7}$ Ha
against a diagonal that contains $10^5$ Ha entries requires an
effective energy precision of $\sim 10^{-12}$ in relative terms,
which is well beyond the scope of VQE or QPE at current shot counts. The
finite-size shift is not only below VQE precision but below the
float$64$ subtraction floor of the nuclear diagonal, and must be
computed analytically as a perturbative shift rather than by direct
diagonalization.

The practical consequence is that the composed nuclear--electronic
Hamiltonian should be solved in block-partitioned form: the nuclear
block is solved first (using the deuteron Hamiltonian in isolation),
the electronic block is solved second (at fixed proton charge
distribution), and the cross-block couplings are treated
perturbatively as classical corrections. This is the quantum-chemistry
analogue of the Born--Oppenheimer approximation, and it is appropriate
whenever the block energies span orders of magnitude. The composed
Hamiltonian in Eq.~\eqref{eq:composed-ne} is most useful as an
\emph{architectural} object: it lets one write down the full
Hamiltonian on a single qubit register, verify correctness by
comparison with the block decomposition, and reason about the cost
of adding cross-block couplings, even if the single-pass solution is
not practical.

\subsection{Known limitation}

The Pauli encoding of the off-diagonal one-body terms in the deuteron
builder of Sec.~\ref{sec:deuteron} carries a known sign issue: the
$a^\dagger_i a_j$ operator for $i \neq j$ is currently encoded using
$XY - YX$ with a real coefficient, which yields the Hermitian
combination $A \otimes A - C \otimes C$ rather than the expected
$A \otimes A + C \otimes C$. This affects only the off-diagonal piece
of the single-particle Hamiltonian, which in the HO basis is
identically zero for the diagonal kinetic term, so the deuteron FCI
spectrum is unaffected. The cross-register hyperfine term in the
composed Hamiltonian is built via the matrix-form code path in
\texttt{nuclear\_electronic.py}, which bypasses the Pauli builder and
yields the correct spectrum directly. We report this as a known issue;
it does not affect the Pauli counts or the $1$-norms reported in
Tables~\ref{tab:deuteron}--\ref{tab:ne-scales}, which are read from
the term list of the encoded Hamiltonian.

%% =====================================================================
\section{Discussion}
\label{sec:discussion}
%% =====================================================================

The results of Sections~\ref{sec:shell}--\ref{sec:ne} suggest a natural
partition of the GeoVac framework into a universal component and a
Coulomb-specific component:

\begin{itemize}
  \item \textbf{Universal (applies to HO, Woods--Saxon, NN potentials).}
    The angular selection rules from Wigner $3j$ symbols and the Gaunt
    integral; the $O(Q^4)$ single-center and $O(Q^{2.5})$ composed
    Pauli scaling; the block structure of the two-body tensor; the
    qubit encoding via Jordan--Wigner on species-specific registers.
  \item \textbf{Coulomb-specific.} The $S^3$ conformal projection
    $p_0 = \sqrt{-2E_n}$; the SO(4) hidden symmetry and the Runge--Lenz
    vector; the $-1/16$ kinetic constant; the Hopf bundle structure
    underlying the fine-structure constant
    conjecture~\cite{Paper2}.
\end{itemize}

The nuclear shell model fits squarely in the universal category. The
Fock rigidity theorem makes this precise: the conformal projection
simply does not survive the transition from Coulomb to HO or
Woods--Saxon, so any claim that the GeoVac $S^3$ machinery extends
directly to nuclei would be incorrect. What does extend is the
angular machinery --- and the results of Sections~\ref{sec:deuteron}
and~\ref{sec:he4} demonstrate that this machinery is sufficient to
produce structurally sparse qubit Hamiltonians for simple nuclei.

\emph{Directions for follow-up work.} (i) Larger
$N_{\mathrm{shells}}$ for converged binding energies; the interesting
physics question is whether the sub-linear Pauli scaling persists as
the basis grows, or whether the Minnesota potential's short range
eventually fills in the angular zeros. (ii) Three-body forces (chiral
EFT at N$^2$LO); these introduce new classes of Pauli terms beyond
the two-body pattern studied here, and their effect on sparsity is an
open question. (iii) Larger nuclei such as $^6$Li, $^{12}$C, and
$^{16}$O; the composed architecture should apply, with nucleon
blocks labelled by single-particle $(n_r, \ell, j)$ orbitals. (iv)
Hardware implementation and comparison with NCSM benchmarks on
actual quantum simulators; the $16$-qubit deuteron is within reach
of current ion-trap and superconducting platforms.

A final note on possible Coulomb-free projections. The Bargmann--Segal
transform~\cite{Bargmann1961} gives a natural holomorphic
representation of the HO Hilbert space, in which the HO Hamiltonian
becomes a first-order differential operator. Whether the
Bargmann--Segal space admits a discrete graph structure whose
Laplacian converges to the HO Hamiltonian --- a strict HO analogue of
the $S^3$ Fock projection --- is an open question that we leave for
future work.

%% =====================================================================
\section{Conclusion}
\label{sec:conclusion}
%% =====================================================================

We have constructed qubit Hamiltonians for the nuclear shell model on
the GeoVac hyperspherical lattice. The HO shell closures and
Mayer--Jensen magic numbers are reproduced from the graph spectrum with
standard spin--orbit and orbital-correction parameters. The deuteron
and He-4 qubit Hamiltonians, built from the Minnesota potential via
Moshinsky--Talmi brackets, show sub-linear Pauli scaling: He-4 requires
only $20\%$ more Pauli terms than the deuteron despite a $12.25\times$
larger Hilbert space. The composed nuclear--electronic deuterium
encoding demonstrates that the multi-block architecture extends to
nucleon--electron systems but faces a $\sim 10^{13}$ coefficient ratio that
requires block-partitioned simulation. The Fock projection rigidity
theorem precisely delineates which parts of the GeoVac framework
transfer to non-Coulomb systems (the angular machinery) and which do
not (the $S^3$ conformal projection). The binding energies reported
here at $N_{\mathrm{shells}}=2$ are not converged and should not be
read as nuclear-structure predictions; they are encoding-validation
benchmarks that verify the qubit Hamiltonian reproduces the matrix
representation.

\begin{thebibliography}{99}

\bibitem{Mayer1949}
M.~G. Mayer,
``On Closed Shells in Nuclei. II,''
Phys.~Rev.~\textbf{75}, 1969 (1949).

\bibitem{Jensen1949}
O.~Haxel, J.~H.~D.~Jensen, and H.~E.~Suess,
``On the ``Magic Numbers'' in Nuclear Structure,''
Phys.~Rev.~\textbf{75}, 1766 (1949).

\bibitem{Suhonen2007}
J.~Suhonen,
\emph{From Nucleons to Nucleus: Concepts of Microscopic Nuclear Theory}
(Springer, Berlin, 2007).

\bibitem{Talmi1993}
I.~Talmi,
\emph{Simple Models of Complex Nuclei: The Shell Model and the
Interacting Boson Model}
(Harwood, Chur, 1993).

\bibitem{Thompson1977}
D.~R. Thompson, M.~Lemere, and Y.~C. Tang,
``Systematic Investigation of Scattering Problems with the
Resonating-Group Method,''
Nucl.~Phys.~A~\textbf{286}, 53 (1977).

\bibitem{Moshinsky1959}
M.~Moshinsky,
``Transformation Brackets for Harmonic Oscillator Functions,''
Nucl.~Phys.~\textbf{13}, 104 (1959).

\bibitem{Talmi1952}
I.~Talmi,
``Nuclear Spectroscopy with Harmonic Oscillator Wavefunctions,''
Helv.~Phys.~Acta \textbf{25}, 185 (1952).

\bibitem{Roggero2020}
A.~Roggero, A.~C.~Y. Li, J.~Carlson, R.~Gupta, and G.~N. Perdue,
``Quantum Computing for Neutrino-Nucleus Scattering,''
Phys.~Rev.~D~\textbf{101}, 074038 (2020).

\bibitem{Hagen2014}
G.~Hagen, T.~Papenbrock, M.~Hjorth-Jensen, and D.~J. Dean,
``Coupled-Cluster Computations of Atomic Nuclei,''
Rep.~Prog.~Phys.~\textbf{77}, 096302 (2014).

\bibitem{Romero2022}
A.~M. Romero \emph{et al.},
``Solving Nuclear Structure Problems with the Adaptive Variational
Quantum Algorithm,''
Phys.~Rev.~C~\textbf{105}, 064317 (2022).

\bibitem{Navratil2000}
P.~Navr\'atil, J.~P. Vary, and B.~R. Barrett,
``Large-Basis Ab Initio No-Core Shell Model and Its Application to
$^{12}$C,''
Phys.~Rev.~Lett.~\textbf{84}, 5728 (2000).

\bibitem{BarrettNavratilVary2013}
B.~R. Barrett, P.~Navr\'atil, and J.~P. Vary,
``Ab Initio No Core Shell Model,''
Prog.~Part.~Nucl.~Phys.~\textbf{69}, 131 (2013).

\bibitem{Fock1935}
V.~Fock,
``Zur Theorie des Wasserstoffatoms,''
Z.~Phys.~\textbf{98}, 145 (1935).

\bibitem{Bander1966}
M.~Bander and C.~Itzykson,
``Group Theory and the Hydrogen Atom,''
Rev.~Mod.~Phys.~\textbf{38}, 330 (1966).

\bibitem{Bargmann1961}
V.~Bargmann,
``On a Hilbert Space of Analytic Functions and an Associated Integral
Transform,''
Commun.~Pure Appl.~Math.~\textbf{14}, 187 (1961).

\bibitem{Foldy1958}
L.~L. Foldy,
``Neutron-Electron Interaction,''
Rev.~Mod.~Phys.~\textbf{30}, 471 (1958).

\bibitem{Friar1979}
J.~L. Friar,
``Nuclear Finite-Size Effects in Light Muonic Atoms,''
Ann.~Phys.~(N.Y.) \textbf{122}, 151 (1979).

\bibitem{Paper0}
J.~Loutey, ``GeoVac Paper 0: The Geometric Packing Construction,''
technical report, Zenodo (2026).

\bibitem{Paper2}
J.~Loutey, ``GeoVac Paper 2: Fine Structure Constant from the Hopf
Bundle,'' technical report, Zenodo (2026). (Conjectural.)

\bibitem{Paper7}
J.~Loutey, ``GeoVac Paper 7: The Dimensionless Vacuum,''
technical report, Zenodo (2026).

\bibitem{Paper14}
J.~Loutey, ``GeoVac Paper 14: Structurally Sparse Qubit Hamiltonians,''
technical report, Zenodo (2026).

\bibitem{Paper17}
J.~Loutey, ``GeoVac Paper 17: Composed Natural Geometries for
Molecules,'' technical report, Zenodo (2026).

\bibitem{Paper22}
J.~Loutey, ``GeoVac Paper 22: Angular Sparsity Theorem,''
technical report, Zenodo (2026).

\end{thebibliography}

\end{document}
